\section{Experiments}
\label{sec:experiments}

We evaluate whether the proposed row-graph models can transfer to a
relational database that is held out from every parameter update. The
experiments address three questions. First, can a small
graph model achieve competitive leave-database-out classification accuracy?
Second, how much of its accuracy remains when the query entity's contextual
target values are removed from the sampled context? Third, what is gained by
conditioning row-level messages on the auxiliary relation graph? Regression
is deliberately left outside the present empirical scope and is reserved for
Section~\ref{sec:exp:regression-placeholder}.

\subsection{Experimental Setup}
\label{sec:exp:setup}

\subsubsection{Benchmarks and Task Selection}
\label{sec:exp:tasks}

We use binary forecasting tasks from RelBench~\citep{robinson2024relbench},
which supplies temporally split prediction tasks over heterogeneous
relational databases. Table~\ref{tab:exp:tasks} lists the ten tasks included
in this initial classification study. They span six databases and markedly
different domains, target relations, and test-set sizes. We define the
evaluation set a priori as the complete ten-task zero-shot classification
suite reported by the Relational Transformer (RT)~\citep{ranjan2026rt}; all
ten tasks are included irrespective of their outcome. Two newer
\texttt{rel-event} classification tasks present in the current codebase lie
outside that published comparison suite. They are excluded from every
aggregate below and will be added in a future expanded comparison once all
methods have been evaluated on them.

Let
\begin{equation}
  \mathfrak T_{\mathrm{clf}}^{\mathrm{eval}}
  \subset \mathfrak T,
  \qquad
  \left|\mathfrak T_{\mathrm{clf}}^{\mathrm{eval}}\right|=10,
\label{eq:exp:evaluation-task-set}
\end{equation}
denote this fixed evaluation set. We retain the task symbol \(\kappa\) from
Section~\ref{sec:method:problem-overview}; \(g\) continues to denote an
individual query rather than a database or task.

\begin{table}[t]
  \centering
  \caption{The ten RelBench classification tasks used in the current
  evaluation. \(N_{\mathrm{test}}\) is the number of queries in the full
  test split. The target column is listed to make explicit which value is
  masked at the seed row and which same-entity values are affected by the
  no-self-label protocol.}
  \label{tab:exp:tasks}
  \footnotesize
  \setlength{\tabcolsep}{4pt}
  \begin{tabular}{llllr}
    \toprule
    Database & Task & Prediction entity & Target column &
    \(N_{\mathrm{test}}\) \\
    \midrule
    \texttt{rel-amazon} & Item Churn & item & \texttt{churn} & 166,842 \\
    \texttt{rel-amazon} & User Churn & user & \texttt{churn} & 351,885 \\
    \texttt{rel-avito} & User Clicks & user & \texttt{num\_click} & 47,996 \\
    \texttt{rel-avito} & User Visits & user & \texttt{num\_click} & 36,129 \\
    \texttt{rel-f1} & Driver DNF & driver & \texttt{did\_not\_finish} & 702 \\
    \texttt{rel-f1} & Driver Top-3 & driver & \texttt{qualifying} & 726 \\
    \texttt{rel-hm} & User Churn & customer & \texttt{churn} & 74,575 \\
    \texttt{rel-stack} & User Badge & user & \texttt{WillGetBadge} & 255,360 \\
    \texttt{rel-stack} & User Engagement & user & \texttt{contribution} & 88,137 \\
    \texttt{rel-trial} & Study Outcome & study & \texttt{outcome} & 825 \\
    \bottomrule
  \end{tabular}
\end{table}

\subsubsection{Compared Methods}
\label{sec:exp:methods}

We compare the following three systems.

\paragraph{Relational Transformer (RT).}
RT~\citep{ranjan2026rt} is the external zero-shot comparator. It serializes
cells from a sampled relational context and processes them with 12
relational-attention blocks. We use the authors' published Standard-test
numbers when comparing to prior work. For the paired robustness experiment,
we evaluate the released leave-database-out checkpoints locally under both
Standard and no-self-label sampling. These two RT result sources are kept
separate throughout the section.

\paragraph{RowGraph-Base.}
This is the base case defined in
Section~\ref{sec:method:base-special-case}. It applies six typed
message-passing layers to the candidate-augmented row graph and sets every
FK condition vector to zero. It therefore tests whether row-level graph
construction, heterogeneous cell encoding, candidate nodes, and the shared
readout already provide transferable inductive bias without the auxiliary
relation graph.

\paragraph{RelGraph.}
This is the complete relation-conditioned model. It first propagates for
two layers over the auxiliary relation graph, uses the resulting relation
representations to construct query-specific FK condition vectors, and then
applies the same six-layer row-graph backbone as RowGraph-Base. RelGraph is
the proposed full model, whereas RowGraph-Base is its principal architectural
ablation.

Table~\ref{tab:exp:model-configurations} summarizes the configurations. Its
counts are trainable network parameters and exclude the frozen MiniLM text
encoder used to precompute cell features. Both
row-graph models use hidden width \(d=256\), precomputed text width
\(d_{\mathrm{text}}=384\), and the same 2048-cell context budget. RelGraph
adds only the two relation-graph layers and condition projections. Thus, the
comparison between the two proposed variants isolates relation conditioning
while holding the sampler, cell encoder, row graph, row-graph depth, and
readout fixed. RT also uses width 256 and text width 384, but its published
implementation uses a 1024-cell context. Consequently, parameter count is a
valid model-size comparison, whereas the current experiment is not a
controlled comparison of FLOPs, memory, or latency.

\begin{table}[t]
  \centering
  \caption{Model configurations used in the classification experiments.
  \(P\) denotes auxiliary relation-graph layers; \(L\)/blocks denotes
  row-message-passing layers for our models or attention blocks for RT.
  RT checkpoints are released pretrained checkpoints, so an optimization
  step count is not applicable to our local evaluation.}
  \label{tab:exp:model-configurations}
  \footnotesize
  \setlength{\tabcolsep}{4pt}
  \begin{tabular}{lrrrrrr}
    \toprule
    Model & Parameters & \(d\) & Cells & \(P\) &
    \(L\)/blocks & Optimization \\
    \midrule
    RT & 22.3M & 256 & 1024 & -- & 12 & released checkpoint \\
    RowGraph-Base & 2.05M & 256 & 2048 & 0 & 6 & 100,000 \\
    RelGraph & 2.78M & 256 & 2048 & 2 & 6 & 100,000 \\
    \bottomrule
  \end{tabular}
\end{table}

\subsubsection{Leave-Database-Out Training and Zero-Shot Inference}
\label{sec:exp:zero-shot-protocol}

For a target database \(D^\star\), let \(\operatorname{db}(\kappa)\)
return the database associated with task \(\kappa\). The source-task pool
used for parameter updates is
\begin{equation}
  \mathfrak T_{\mathrm{train}}(D^\star)
  =
  \left\{
    \kappa\in\mathfrak T
    \,\middle|\,
    \operatorname{db}(\kappa)\neq D^\star
  \right\}.
\label{eq:exp:leave-db-out-set}
\end{equation}
No query from \(D^\star\) contributes to an optimization gradient. A
separate randomly initialized leave-database-out run is trained for each
reported target task, including tasks that share the same held-out
database. The source stream is the implementation's complete set of tasks
from the remaining databases; it contains both binary and continuous-label
source tasks. The present section evaluates only binary target tasks. We
state this mixed-source detail explicitly because the eventual regression
methodology must document the continuous-label source objective before the
paper is fully reproducible.

We call this protocol \emph{leave-database-out zero-shot transfer}: after
source training, the learned parameters are applied to test queries from
\(D^\star\) without task- or database-specific fine-tuning. ``Zero-shot''
here constrains parameter learning; it does not mean that the target
database is absent at inference. Its schema, rows, timestamps, and FK
structure are the input from which each query graph is constructed. This
usage is analogous to inductive zero-shot graph inference, where a fixed
model is evaluated on an unseen graph while that graph's structure remains
available as input~\citep{galkin2024ultra}.

There is one important qualification. Target-task validation labels are
used to select a checkpoint, although they are never used to update model
parameters. Our precise protocol is therefore
\emph{leave-database-out training with target-validation checkpoint
selection}, followed by zero-shot test inference. It is stronger than
target-database fine-tuning because the target database contributes no
gradient, but weaker than a source-only selection protocol in which no
target label is consulted before test. We report this distinction rather
than describing the target database as unseen in every possible sense.

\subsubsection{Standard and No-Self-Label Evaluation}
\label{sec:exp:nsl-protocol}

In both evaluation regimes, the seed target cell \(n_g^\star\) is masked as
specified in Section~\ref{sec:method:problem-overview}. The regimes differ
only in whether the sampler may expose contextual target-column cells from
other rows associated with the same query entity.

To define this precisely, let \(\mathcal C^{\mathrm{DB}}\) denote the set of
database cells and let \(\mathcal P(r)\) be the set of database rows
referenced by outgoing FK links of row \(r\). A sampled non-seed row \(r\)
shares a query entity with the seed whenever
\(\mathcal P(r)\cap\mathcal P(r_g^\star)\neq\varnothing\). The set of
copyable self-label cells is
\begin{equation}
  \mathcal C_g^{\mathrm{self}}
  =
  \left\{
    n\in\mathcal C^{\mathrm{DB}}
    \,\middle|\,
    \begin{array}{l}
      r(n)\neq r_g^\star,\quad
      \tblname(r(n))=\tblname(r_g^\star),\\[-1mm]
      c(n)=c_g^\star,\quad
      \mathcal P(r(n))\cap\mathcal P(r_g^\star)\neq\varnothing
    \end{array}
  \right\},
\label{eq:exp:self-label-set}
\end{equation}
where \(r(n)\) and \(c(n)\) are the row and column of cell \(n\).
Operationally, these are contextual values of the same target column on
other task-table rows that share at least one FK parent with the seed.

The \emph{Standard} regime sets the sampler's self-label dropout to zero,
so a cell in \(\mathcal C_g^{\mathrm{self}}\) remains visible if ordinary
context sampling selects it. The \emph{no-self-label} (NSL) regime sets
that dropout to one, so every selected cell in
\(\mathcal C_g^{\mathrm{self}}\) is skipped. NSL does \emph{not} remove
other features of those rows, target-column cells associated with different
entities, or explicitly delete FK edges. The final sampled graph can still
change when replacement cells enter the fixed budget. It therefore measures
dependence on one specific shortcut---the query entity's contextual target
values---rather than removing all observed labels or all relational signal. Standard and
NSL always evaluate the same model checkpoint on the same test queries.

The deletion occurs inside serialization under a fixed cell budget. A
skipped self-label cell can therefore free a position for a later context
cell. NSL is consequently a resampled-context intervention rather than a
strict operation that zeros coordinates in an otherwise identical graph.
This detail matters when interpreting the few cases in which NSL exceeds
Standard: such an increase does not imply that deleting information is
intrinsically beneficial, because the replacement context also changed.

\subsubsection{Training, Model Selection, and Metrics}
\label{sec:exp:optimization}

Both proposed variants are trained from random initialization for 100,000
optimization steps with batch size 32. We use AdamW with maximum learning
rate \(10^{-3}\), weight decay \(0.1\), gradient-norm clipping at \(1.0\),
and a linear OneCycle schedule with warm-up fraction \(0.2\). The sampler
uses a 2048-cell budget, reverse-FK width \(B_{\max}=256\), and no
train-time self-label dropout. Text cells use fixed 384-dimensional features
precomputed by \texttt{all-MiniLM-L12-v2}; the text encoder receives no
gradient and is not included in the parameter counts. All reported runs use
random seed zero and no warm-start checkpoint.

At each validation event, a checkpoint is scored on both regimes. Let
\(A_{\mathrm{val},\kappa}^{\mathrm{Std}}(u)\) and
\(A_{\mathrm{val},\kappa}^{\mathrm{NSL}}(u)\) be validation AUROC for
target task \(\kappa\) at training step \(u\). We choose
\begin{equation}
  u_\kappa^\star
  =
  \underset{u}{\arg\max}\;
  \left[
    A_{\mathrm{val},\kappa}^{\mathrm{Std}}(u)
    +
    A_{\mathrm{val},\kappa}^{\mathrm{NSL}}(u)
  \right].
\label{eq:exp:checkpoint-selection}
\end{equation}
Each validation regime is capped at 64 batches, or approximately 2048
queries. The selected checkpoint is then evaluated on every query in the
full test split under both regimes. Test results do not enter
Eq.~\eqref{eq:exp:checkpoint-selection}.

For model \(M\) and evaluation task \(\kappa\), define the full-test scores
\begin{equation}
  A_{M,\kappa}^{\mathrm{Std}}
  =100\times\mathrm{AUROC}_{M,\kappa}^{\mathrm{Std}},
  \qquad
  A_{M,\kappa}^{\mathrm{NSL}}
  =100\times\mathrm{AUROC}_{M,\kappa}^{\mathrm{NSL}}.
\label{eq:exp:task-auroc}
\end{equation}
Hence, every table entry is in AUROC points. We aggregate tasks with equal
weight, not in proportion to their test-set sizes:
\begin{align}
  \overline A_M^{\mathrm{Std}}
  & =
  \frac{1}{|\mathfrak T_{\mathrm{clf}}^{\mathrm{eval}}|}
  \sum_{\kappa\in\mathfrak T_{\mathrm{clf}}^{\mathrm{eval}}}
  A_{M,\kappa}^{\mathrm{Std}},
  \label{eq:exp:macro-standard}\\
  \overline A_M^{\mathrm{NSL}}
  & =
  \frac{1}{|\mathfrak T_{\mathrm{clf}}^{\mathrm{eval}}|}
  \sum_{\kappa\in\mathfrak T_{\mathrm{clf}}^{\mathrm{eval}}}
  A_{M,\kappa}^{\mathrm{NSL}}.
  \label{eq:exp:macro-nsl}
\end{align}
For paired runs, sensitivity to the NSL intervention is summarized by
\begin{equation}
  \operatorname{Drop}_{\mathrm{NSL}}(M)
  =
  \overline A_M^{\mathrm{Std}}
  -
  \overline A_M^{\mathrm{NSL}}.
\label{eq:exp:self-label-drop}
\end{equation}
A smaller value means that less aggregate performance is lost when the
self-label cells in Eq.~\eqref{eq:exp:self-label-set} are unavailable. It
must be interpreted together with the absolute NSL AUROC and does not by
itself prove a particular reasoning mechanism.

\subsection{Parameter-Efficient Leave-Database-Out Zero-Shot Classification}
\label{sec:exp:standard-results}

\subsubsection{Overall Standard Performance}
\label{sec:exp:overall-standard}

Table~\ref{tab:exp:standard-results} reports Standard full-test AUROC. The
RT column uses the values published in its zero-shot classification
table~\citep{ranjan2026rt}; the two graph-model columns use our complete
full-test evaluations. The latter follow the target-validation checkpoint
selection in Eq.~\eqref{eq:exp:checkpoint-selection}; accordingly, this is
a comparison under our declared leave-database-out protocol, not under
source-only model selection. RelGraph reaches a macro-average of \(69.7\), equal
to RT's published \(69.7\) after rounding to one decimal place, and improves
over RowGraph-Base by \(2.1\) points. Before rounding, the macro averages are
\(69.71\) for published RT, \(67.574\) for RowGraph-Base, and \(69.663\)
for RelGraph. We therefore describe RelGraph as matching the reported RT
macro-average at the displayed precision, not as statistically
indistinguishable from RT.

The average conceals substantial task variation. RelGraph exceeds the
published RT score on six of ten tasks: both Avito tasks, Driver Top-3,
H\&M User Churn, Stack User Engagement, and Trial Study Outcome. RT remains
stronger on both Amazon tasks, Driver DNF, and Stack User Badge.
RowGraph-Base obtains the best score on H\&M User Churn and Trial Study
Outcome, showing that auxiliary relation conditioning is not uniformly
beneficial. Trial Study Outcome is also a clear failure case for all three
systems: scores close to \(50\) indicate that none extracts strong ranking
signal on this task.

\begin{table}[t]
  \centering
  \caption{Standard-test classification performance
  (\(100\times\mathrm{AUROC}\), higher is better). RT values are the
  published zero-shot results; RowGraph-Base and RelGraph values are full
  test-set evaluations of the selected checkpoints. The macro average gives
  equal weight to each of the ten tasks. Bold marks the largest displayed
  value in each row.}
  \label{tab:exp:standard-results}
  \small
  \begin{tabular}{lrrr}
    \toprule
    Task & RT & RowGraph-Base & RelGraph \\
    \midrule
    Amazon / Item Churn & \textbf{70.9} & 68.4 & 66.9 \\
    Amazon / User Churn & \textbf{64.0} & 63.1 & 63.7 \\
    Avito / User Clicks & 59.5 & 53.8 & \textbf{61.1} \\
    Avito / User Visits & 61.8 & 61.2 & \textbf{62.4} \\
    F1 / Driver DNF & \textbf{81.2} & 75.3 & 79.7 \\
    F1 / Driver Top-3 & 89.3 & 88.0 & \textbf{89.6} \\
    H\&M / User Churn & 62.8 & \textbf{67.2} & 66.8 \\
    Stack / User Badge & \textbf{80.1} & 75.1 & 76.8 \\
    Stack / User Engagement & 75.7 & 70.0 & \textbf{76.3} \\
    Trial / Study Outcome & 51.8 & \textbf{53.6} & 53.2 \\
    \midrule
    Macro average & \textbf{69.7} & 67.6 & \textbf{69.7} \\
    \bottomrule
  \end{tabular}
\end{table}

\subsubsection{Performance Relative to Model Size}
\label{sec:exp:parameter-efficiency}

RelGraph contains \(2.78\)M trainable network parameters, compared with
\(22.3\)M for RT; neither count includes the frozen text encoder used to
precompute cell features. Its relative parameter count and reduction are
\begin{equation}
  \frac{2.78}{22.3}=0.125,
  \qquad
  1-\frac{2.78}{22.3}=0.875.
\label{eq:exp:parameter-ratio}
\end{equation}
Thus, RelGraph uses approximately one eighth as many parameters, an
\(87.5\%\) reduction, while matching RT's published Standard macro-AUROC
after rounding. RowGraph-Base is smaller still at \(2.05\)M parameters and
trails RT by only \(2.1\) macro points. These results provide evidence that
competitive leave-database-out transfer on this ten-task suite does not
require a large Transformer backbone: substantial cross-database accuracy
can be encoded by a compact row-level graph model with parameters shared
across schemas and tasks within each leave-database-out run.

This is specifically a claim about \emph{parameter efficiency}. Because RT
and the proposed models use different context budgets and operators,
Table~\ref{tab:exp:model-configurations} does not establish lower inference
time, lower peak memory, or fewer floating-point operations. Those quantities
require direct hardware-controlled measurement.

\subsection{Robustness to Same-Entity Target Values}
\label{sec:exp:nsl-results}

\subsubsection{Aggregate Robustness}
\label{sec:exp:nsl-aggregate}

Table~\ref{tab:exp:nsl-results} gives paired Standard and NSL evaluations.
For RT, both columns in this table come from our local evaluation of the
same released checkpoint; they are not formed by pairing a published
Standard score with a locally reproduced NSL score. This distinction makes
the within-model drop in Eq.~\eqref{eq:exp:self-label-drop} meaningful.

Under NSL, RelGraph retains a macro-average AUROC of \(62.3\), compared with
\(52.1\) for RT, a gain of \(10.2\) points at the reported precision. Its
paired Standard-to-NSL drop is \(7.3\) points, much
smaller than RT's \(17.8\)-point drop. RowGraph-Base is the strongest NSL
variant in aggregate, reaching \(63.9\) with a drop of only \(3.7\) points.
The two proposed row-graph models therefore retain substantially more
ranking performance than RT after removal of same-entity contextual target
values under their respective checkpoint-selection procedures. In
particular, our row-graph checkpoint criterion explicitly includes NSL
validation performance, whereas the released RT checkpoints were not
selected with Eq.~\eqref{eq:exp:checkpoint-selection}. The comparison
therefore measures the final protocol-level robustness of the available
systems; it does not isolate architecture while controlling model selection.

This result does not imply that the row-graph models never use observed
labels, nor that the remaining prediction is purely structural. NSL keeps
other entities' target values, every non-target feature, timestamps, and the
sampled FK neighborhood. The supported conclusion is narrower: the average
leave-database-out performance of the row-graph models is less dependent on the
specific copyable self-label set in
Eq.~\eqref{eq:exp:self-label-set}, and useful predictive signal remains in
the rest of the relational context.

\begin{table}[t]
  \centering
  \caption{Paired Standard and no-self-label (NSL) performance
  (\(100\times\mathrm{AUROC}\)). Every pair evaluates one checkpoint on
  the same test queries. RT numbers here are our local paired reproduction
  using released checkpoints; the published RT values used for prior-work
  comparison appear separately in Table~\ref{tab:exp:standard-results}.
  Bold marks the best NSL score in each task row and the smallest aggregate
  NSL drop.}
  \label{tab:exp:nsl-results}
  \footnotesize
  \setlength{\tabcolsep}{4pt}
  \begin{tabular}{lrrrrrr}
    \toprule
    & \multicolumn{2}{c}{RT, reproduced}
    & \multicolumn{2}{c}{RowGraph-Base}
    & \multicolumn{2}{c}{RelGraph} \\
    \cmidrule(lr){2-3}\cmidrule(lr){4-5}\cmidrule(lr){6-7}
    Task & Standard & NSL & Standard & NSL & Standard & NSL \\
    \midrule
    Amazon / Item Churn
      & 70.2 & 51.8 & 68.4 & \textbf{71.1} & 66.9 & 65.1 \\
    Amazon / User Churn
      & 64.6 & 54.9 & 63.1 & 61.5 & 63.7 & \textbf{61.8} \\
    Avito / User Clicks
      & 58.5 & 53.9 & 53.8 & 59.5 & 61.1 & \textbf{61.2} \\
    Avito / User Visits
      & 62.1 & 48.0 & 61.2 & 61.2 & 62.4 & \textbf{62.1} \\
    F1 / Driver DNF
      & 79.1 & 69.7 & 75.3 & 46.4 & 79.7 & \textbf{70.8} \\
    F1 / Driver Top-3
      & 89.1 & 46.3 & 88.0 & \textbf{87.1} & 89.6 & 67.0 \\
    H\&M / User Churn
      & 64.2 & 39.6 & 67.2 & 61.1 & 66.8 & \textbf{63.6} \\
    Stack / User Badge
      & 80.6 & 54.3 & 75.1 & \textbf{73.1} & 76.8 & 65.6 \\
    Stack / User Engagement
      & 78.3 & 49.9 & 70.0 & \textbf{64.0} & 76.3 & 52.9 \\
    Trial / Study Outcome
      & 52.4 & 52.4 & 53.6 & \textbf{53.6} & 53.2 & 53.2 \\
    \midrule
    Macro average
      & 69.9 & 52.1 & 67.6 & \textbf{63.9} & 69.7 & 62.3 \\
    NSL drop \(\downarrow\)
      & \multicolumn{2}{c}{17.8}
      & \multicolumn{2}{c}{\textbf{3.7}}
      & \multicolumn{2}{c}{7.3} \\
    \bottomrule
  \end{tabular}
\end{table}

\subsubsection{Task-Level Behavior and a Natural Control}
\label{sec:exp:nsl-task-level}

The aggregate advantage is not produced by one isolated task. RelGraph's
NSL score exceeds RT's reproduced NSL score on all ten tasks, while
RowGraph-Base and RelGraph split the task-level NSL wins between them. The
largest RT drops occur on Driver Top-3, H\&M User Churn, Stack User Badge,
and Stack User Engagement. In contrast, several row-graph scores change
only modestly after self-label removal.

Some row-graph tasks improve under NSL, most notably RowGraph-Base on Amazon
Item Churn and Avito User Clicks. As noted in
Section~\ref{sec:exp:nsl-protocol}, this should not be interpreted as causal
evidence that withholding a target value creates information. Removing a
cell during fixed-budget serialization permits another cell later in the
sample to enter, so the contextual graph is not held perfectly fixed.

Trial Study Outcome provides a useful natural control. Its sampled contexts
contain no cells satisfying Eq.~\eqref{eq:exp:self-label-set}; consequently,
Standard and NSL are equal up to numerical rounding for each locally
evaluated model. This is consistent with the NSL intervention leaving an
unaffected task unchanged. At the same time, all three AUROCs remain close
to chance, emphasizing that robustness to self-label removal is not by
itself sufficient for accurate relational prediction.

\subsection{Ablation of Relation Conditioning}
\label{sec:exp:relation-ablation}

The difference between RowGraph-Base and RelGraph isolates the auxiliary
relation graph and its FK condition vectors. Adding relation conditioning
increases the parameter count from \(2.05\)M to \(2.78\)M and improves
Standard macro-AUROC from \(67.6\) to \(69.7\), a gain of \(2.1\) points.
The gains are especially visible for Avito User Clicks
(\(+7.3\) points), Driver DNF (\(+4.5\)), and Stack User Engagement
(\(+6.3\)). These tasks show that query-specific relation information can
make a compact row-graph backbone more competitive under the ordinary
evaluation distribution.

Relation conditioning is not the sole source of NSL robustness. The NSL
macro-average decreases from \(63.9\) for RowGraph-Base to \(62.3\) for
RelGraph, and the effect is heterogeneous across tasks. RelGraph improves
NSL by \(24.4\) points on Driver DNF and by \(2.5\) on H\&M User Churn, but
is \(20.1\) points lower on Driver Top-3 and \(11.0\) points lower on Stack
User Engagement. Thus, the present ablation supports two distinct
conclusions: relation conditioning improves average Standard accuracy, while
the base row-graph construction already supplies much of the resistance to
same-entity target removal. Understanding why relation conditioning
helps some NSL tasks and harms others is an open target for controlled
multi-seed analysis.

\subsection{Limitations and Current Empirical Scope}
\label{sec:exp:limitations}

The current results establish a detailed first comparison, but they do not
close the empirical study.

\paragraph{Single training seed.}
RowGraph-Base and RelGraph are currently reported for seed zero. We
therefore make no claim of statistical significance, and matching after
one-decimal rounding should not be confused with an equivalence test.
Multi-seed training is required to quantify optimization variance,
especially on the small F1 and Trial test sets.

\paragraph{Target-validation selection.}
No target-database query contributes a gradient, but
Eq.~\eqref{eq:exp:checkpoint-selection} uses target validation labels and
explicitly rewards both Standard and NSL validation performance. The
reported NSL robustness is therefore partly a model-selection objective,
not a fully post hoc property. A stricter future experiment should choose
the checkpoint using source validation only and evaluate the target
database exactly once.

\paragraph{Initial ten-task suite.}
The study currently covers the fixed ten-task RT comparison suite described
in Section~\ref{sec:exp:tasks}. It does not yet include the two newer
\texttt{rel-event} classification tasks or additional external models.
Conclusions are limited to this declared set, and future tasks must be added
without retroactively changing the aggregates reported here.

\paragraph{Mixed source tasks.}
The training stream contains the available classification and regression
tasks from source databases, even though this section evaluates only
classification. The continuous-label candidate construction and loss must
be specified in the regression extension before the final manuscript is
fully self-contained.

\paragraph{Parameter versus system efficiency.}
The \(87.5\%\) reduction relative to RT concerns trainable parameter count.
Different cell budgets and operators prevent a direct claim about compute,
latency, throughput, or memory. Those comparisons require profiling under a
shared hardware and batching protocol.

\paragraph{Scope of NSL.}
NSL removes only the query entity's contextual target-column cells. It does
not remove labels belonging to other entities, all temporal correlations,
or all relational context, and it can alter which later cells fit into the
budget. Accordingly, NSL diagnoses sensitivity to a particular shortcut;
it does not prove that the retained performance arises from a unique form
of structural reasoning.

\subsection{Extension to Regression}
\label{sec:exp:regression-placeholder}

\emph{Placeholder. This subsection will report regression benchmarks,
metrics, baselines, Standard and no-self-label results, and relation-graph
ablations after the continuous prediction formulation and experimental
protocol are finalized. No regression target result contributes to any
classification average in this section.}
